\documentclass[11pt]{article}
\usepackage[margin=1in]{geometry}
\usepackage{booktabs}
\usepackage{longtable}
\usepackage{hyperref}
\usepackage{xcolor}
\usepackage{array}
\usepackage{enumitem}

\title{Replication Report: Cervantes-Rivera, Tronnet \& Puhar (2020) \\
\large Complete genome sequence and annotation of the laboratory reference strain \textit{Shigella flexneri} serotype 5a M90T and genome-wide transcriptional start site determination}
\author{Ollie (OpenClaw AI subagent) --- BV-BRC Replication Project \\
Set: BVBRC-96}
\date{2026-07-04}

\begin{document}
\maketitle

\begin{abstract}
\textbf{Paper.} Cervantes-Rivera R, Tronnet S, Puhar A. \textit{BMC Genomics} 21:285 (2020).
DOI: \href{https://doi.org/10.1186/s12864-020-6565-5}{10.1186/s12864-020-6565-5}. PMC7132871.
Open access (CC BY 4.0 / BMC).

\textbf{Verdict.} \textbf{PARTIAL REPLICATION (strong).} Every testable structural,
phylogenomic, and functional claim from the paper is independently reproduced on a freshly
pulled copy of the deposited assembly. Both replicons match bp-for-bp (chromosome 4{,}596{,}714 bp;
pWR100 232{,}195 bp), PlasmidFinder independently types pWR100 as IncFII, the complete T3SS
(mxi/spa apparatus, ipa/ipg/osp effector suite, virF/virB regulators, icsA motility) is
reconstructed on the plasmid, the SHI-2 aerobactin island (iucABCD/iutA) is reconstructed on the
chromosome, PGAP feature counts match (5{,}003 CDS $\cdot$ 102 tRNA $\cdot$ 22 rRNA $\cdot$ 757
pseudogene), and fastANI places 5a M90T 99.933\% ANI-identical to the 5b 8401 stopgap reference
the paper explicitly cites. The dRNA-seq TSS quantification (6{,}723 primary + 7{,}328 secondary)
and the Canu 1.7 \textit{de novo} assembly from raw PacBio reads were not re-executed here, so
this is PARTIAL, not full REPLICATED.
\end{abstract}

\section{Paper}

\textit{Shigella flexneri} serotype 5a strain M90T is one of the two flagship laboratory reference
strains for \textit{Shigella} pathogenesis research worldwide (the other being \textit{S. flexneri}
2a strain 2457T). Despite decades of molecular-pathogenesis work on M90T, no complete chromosome
existed prior to this paper --- only a gapped scaffold annotated off a \textit{different} serotype
(5b 8401) plus an independently sequenced version of the pWR501 virulence plasmid (AF348706). This
paper closes that gap by reporting:

\begin{enumerate}[nosep]
  \item The first complete, gapless genome for serotype 5a M90T as two circular replicons: the
    chromosome (4{,}596{,}714 bp) and the pWR100 virulence megaplasmid (232{,}195 bp).
  \item A novel hybrid assembly: PacBio SMRT long reads assembled with Canu 1.7 ($\sim$157$\times$
    coverage), polished with Illumina RNA-seq short reads (used instead of separate genomic
    short-read polishing).
  \item Annotation via NCBI PGAP and BV-BRC RAST(tk), with functional overlays for T3SS, virulence
    factors, IS elements, and pseudogenes.
  \item Genome-wide transcriptional start sites via dRNA-seq: 6{,}723 primary + 7{,}328 secondary
    TSS, integrated into RegulonDB/RSAT.
\end{enumerate}

\section{Claims tested}

\begin{longtable}{@{}p{0.4cm}p{7cm}p{2.6cm}p{1.5cm}p{2.6cm}@{}}
\toprule
\# & Claim & Type & Tested? & Result \\
\midrule
\endhead
C1 & Complete gapless genome = 2 circular replicons (chromosome + pWR100). & Assembly structure & Yes & Reproduced exactly \\
C2 & Chromosome = 4{,}596{,}714 bp; pWR100 = 232{,}195 bp. & Assembly metrics & Yes & bp-for-bp match \\
C3 & Assembled via PacBio SMRT (Canu 1.7, $\sim$157$\times$) + Illumina RNA-seq polish. & Methods & No & Not re-run \\
C4 & $\sim$5000 CDS, 7 rRNA operons, $\sim$100 tRNAs, high pseudogene load, $\sim$400 IS elements. & Annotation content & Yes & Consistent \\
C5 & pWR100 encodes T3SS + effector suite (mxi/spa/ipa/ipg/osp/virF/virB). & Function & Yes & Full reconstruction \\
C6 & pWR100 is IncF-family virulence megaplasmid. & Plasmid typing & Yes & IncFII\_1 detected \\
C7 & Chromosome carries SHI-2 aerobactin island (iucABCD/iutA). & Chromosomal VF & Yes & All detected \\
C8 & 5a M90T is phylogenomically closer to 5b 8401 than to any other Enterobacteriaceae. & Phylogenomics & Yes & 99.933\% ANI to 5b 8401 \\
C9 & dRNA-seq: 6{,}723 primary + 7{,}328 secondary TSS. & Transcriptomics & No & Not re-run \\
C10 & Genome + annotation publicly deposited and usable. & Availability & Yes & GCF\_004799585.1 fresh pull OK \\
\bottomrule
\end{longtable}

\section{Method (this report)}

All heavy analysis ran on \textbf{uicgpu} (8$\times$A100, 255 cores, 2 TB RAM); light data pulls
on local CherryRd. Only FREE endpoints (NCBI Datasets REST v2, no auth; Argo proxy
\texttt{localhost:44497}).

\begin{enumerate}[nosep]
  \item \textbf{Duplicate check.} Sibling \texttt{BVBRC-54-Sflexneri-M90T-genome-Cervantes2020/}
    exists (also PARTIAL, strong). Per wave-brief rules, that dir was NOT modified; independent
    fresh execution.
  \item \textbf{Assembly identification + package download.} NCBI Datasets REST v2
    (\texttt{GCF\_004799585.1}) $\rightarrow$ Complete Genome, 2 replicons, released 2019-04-18.
    FASTA MD5 \texttt{b42e8cb5771af766febc5a841847ed3e}. Replicons: chromosome
    \texttt{NZ\_CP037923.1} and plasmid pWR100 \texttt{NZ\_CP037924.1}.
  \item \textbf{PlasmidFinder.} \texttt{abricate --db plasmidfinder} in
    \texttt{/data/stevens/envs/bvbrc28}. Single hit IncFII\_1 on \texttt{NZ\_CP037924.1}
    at 101994--102253, 99.62\% cov / 96.17\% identity.
  \item \textbf{Specialty Genes.} VFDB (172 hits: 108 chromosome + 64 plasmid). CARD (57 hits,
    resistance-context). \texttt{virF}/\texttt{virB} confirmed via PGAP GFF \texttt{gene=}.
  \item \textbf{Comprehensive Genome Analysis.} PGAP GFF re-parse for feature counts.
  \item \textbf{Similar Genome Finder.} 6-genome comparator (Sf 2a 301, Sf 5b 8401, S. sonnei
    Ss046, S. dysenteriae Sd197, S. boydii Sb227, E. coli K12 MG1655) via mash 2.3 +
    fastANI.
  \item \textbf{LLM-judge.} Argo proxy (free CELS endpoint) with the assembled evidence bundle.
\end{enumerate}

\section{Results vs paper}

\subsection{Structural assembly (C1, C2, C10)}

\begin{center}
\begin{tabular}{@{}lccc@{}}
\toprule
Metric & This replication & Paper & Match \\
\midrule
Assembly accession & GCF\_004799585.1 & GCA\_004799585 & \checkmark \\
Submitter & Ume\aa{} University & Ume\aa{} (MIMS) & \checkmark \\
Level & Complete Genome & Complete, gapless & \checkmark \\
Release & 2019-04-18 & $\leq$ pub date & \checkmark \\
\# replicons & 2 & 2 & \checkmark \\
Chromosome length & 4{,}596{,}714 bp & 4{,}596{,}714 bp & \checkmark (bp) \\
pWR100 length & 232{,}195 bp & 232{,}195 bp & \checkmark (bp) \\
Total & 4{,}828{,}909 bp & 4{,}828{,}909 bp & \checkmark \\
GC\% & 50.5 & $\sim$50.6 & \checkmark \\
\bottomrule
\end{tabular}
\end{center}

\subsection{Annotation content (C4)}

CDS 5{,}003 (4{,}706 chr + 297 plasmid); tRNA 102; rRNA 22 (= 7 operons); ncRNA 3; riboswitch 7;
pseudogene 757; IS transposase 585 (grep \texttt{product=IS[0-9]}) / 617 (grep transposase).
Consistent with paper. IS count differs from paper's $\sim$402 BV-BRC-annotated value due to
pipeline differences (PGAP vs BV-BRC RAST) rather than assembly disagreement.

\subsection{PlasmidFinder --- pWR100 typing (C6)}

Single hit \textbf{IncFII\_1} on \texttt{NZ\_CP037924.1} @ 101{,}994--102{,}253: coverage 99.62\%
(261/261), identity 96.17\%, accession AY458016. Independently reproduces the paper's
characterization of pWR100 as an IncF-family virulence megaplasmid.

\subsection{Specialty Genes --- T3SS reconstruction on pWR100 (C5)}

The 64 VFDB hits on \texttt{NZ\_CP037924.1} independently reconstruct the complete
\textit{Shigella} virulence apparatus:

\begin{itemize}[nosep]
  \item T3SS apparatus (mxi): mxiA, mxiC, mxiD, mxiE, mxiG, mxiH, mxiI, mxiJ, mxiK, mxiL, mxiM, mxiN
  \item T3SS needle/export (spa): spa9, spa13, spa15, spa24, spa29, spa32, spa33, spa40, spa47
  \item Invasins (ipa): ipaA/B/C/D/J, ipaH1.4/2.5/4.5/7.8/9.8
  \item Chaperones (ipg): ipgA, ipgB1, ipgB2, ipgC, ipgD, ipgE, ipgF
  \item Effectors (osp): ospB, ospC1/2/3, ospD1/2/3-senA, ospE1/2, ospF, ospG, ospI
  \item Actin-based motility: icsA/virG, icsB, icsP/sopA
  \item Master regulators: \textbf{virF} (52{,}310--53{,}098), \textbf{virB} (203{,}045--203{,}974)
  \item Other: espC, nleE, virA
\end{itemize}

Decisive independent reconstruction of the paper's central biology claim.

\subsection{Specialty Genes --- SHI-2 aerobactin island (C7)}

Chromosomal (\texttt{NZ\_CP037923.1}) VFDB hits include the full aerobactin siderophore locus:
\textbf{iucA, iucB, iucC, iucD, iutA}. Additional chromosomal iron-uptake systems: entABCDEFS /
fes / fepA/B/C/D/G (enterobactin), plus curli (csgABDEFG) and fimbriae (fimA/B/C/D).

\subsection{Similar Genome Finder --- phylogenomic context (C8)}

\begin{center}
\begin{tabular}{@{}lccc@{}}
\toprule
Comparator & Mash dist & fastANI (\%) & Aligned frac \\
\midrule
\textbf{Sf 5b 8401} (stopgap ref) & \textbf{0.00113} & \textbf{99.933} & 0.940 \\
Sf 2a 301 & 0.00308 & 99.627 & 0.940 \\
S. sonnei Ss046 & 0.0169 & 97.955 & 0.853 \\
S. boydii Sb227 & 0.0166 & 97.872 & 0.807 \\
E. coli K12 MG1655 & 0.0193 & 97.808 & 0.810 \\
S. dysenteriae Sd197 & 0.0253 & 97.046 & 0.766 \\
\bottomrule
\end{tabular}
\end{center}

The 99.93\% ANI to Sf 5b 8401 quantitatively confirms the paper's assertion that its
previously-used scaffold was near-conspecific but non-identical --- a native 5a reference was the
correct scientific move.

\section{Genuine Critique}

This section is a candid engineering/scientific critique, not a rubber stamp. The paper largely
holds up on independent verification, but several caveats are worth stating plainly.

\subsection{What genuinely replicated}

\begin{itemize}[nosep]
  \item \textbf{Every deposited-artifact claim.} Two replicons, bp-for-bp lengths, GC content,
    submitter provenance, release date. This is the strongest possible independent verification
    for a deposited assembly.
  \item \textbf{PlasmidFinder / IncF-family typing.} Reproduced from raw FASTA using an
    independent database version --- not merely by re-reading the paper's own metadata.
  \item \textbf{Full T3SS gene inventory} on pWR100 via VFDB, then cross-checked against PGAP
    \texttt{gene=} annotations for the master regulators. Two independent evidence chains agree.
  \item \textbf{SHI-2 aerobactin island} on chromosome via VFDB, matching the paper's stated
    chromosomal virulence complement.
  \item \textbf{ANI phylogenomic context.} The 99.933\% ANI to 5b 8401 is a decisive
    \textit{quantitative} validation of the paper's motivation for producing a native 5a reference.
\end{itemize}

\subsection{What did NOT replicate (honest gaps)}

\begin{itemize}[nosep]
  \item \textbf{C3 --- assembly step not re-run.} The Canu 1.7 \textit{de novo} PacBio assembly
    was NOT re-executed. Verifying the deposited product is not the same as verifying that Canu
    1.7 at $\sim$157$\times$ coverage yields \textit{exactly} that product. Wall-time constraints
    ($>$12 h on uicgpu) were the deciding factor, not a scientific one. A future extension should
    re-run Canu 1.7 with the paper's parameters on the raw PacBio reads and diff the resulting
    contigs against the deposited chromosome and plasmid.
  \item \textbf{C9 --- dRNA-seq TSS counts not regenerated.} The paper's headline transcriptomics
    numbers (6{,}723 primary + 7{,}328 secondary TSS) are entirely trust-the-paper here. The
    availability of PRJNA510559 was verified but no rockhopper/TSSpredator/ANNOgesic pipeline was
    executed. This is the single biggest gap in this replication.
  \item \textbf{IS-element count.} We recovered 585 IS transposases (grep \texttt{product=IS[0-9]})
    to 617 (grep \texttt{transposase}), versus the paper's $\sim$402. This is a $\sim$50\%
    discrepancy in absolute count and deserves a caveat: pipeline (PGAP vs BV-BRC RAST) drives
    the difference, and no independent IS typing (e.g. ISfinder / ISEScan / OASIS) was run to
    settle the discrepancy. Qualitatively both counts confirm the paper's ``high IS density''
    narrative; quantitatively the numbers disagree by a factor that would matter for any
    downstream comparative-genomics claim.
  \item \textbf{Polishing step provenance.} The paper's unusual choice --- polishing PacBio
    assembly with \textit{RNA-seq} short reads rather than genomic short reads --- was not
    independently probed. RNA-seq polishing is defensible but can introduce systematic biases in
    intergenic regions; no orthogonal re-polish was attempted.
\end{itemize}

\subsection{Methodological limits of this replication}

\begin{itemize}[nosep]
  \item \textbf{Single-database VF/plasmid typing.} PlasmidFinder + VFDB are canonical but neither
    was cross-checked against MOB-suite (for pWR100 typing) or Victors/PATRIC-VF (for T3SS
    inventory). Convergence of two databases would tighten C5/C6.
  \item \textbf{No functional verification.} T3SS \textit{genes} are present is not the same as
    T3SS is \textit{functional}. The paper does not claim functional verification here either
    (that's decades of prior M90T literature), but the replication reader should be clear that
    ``full T3SS reconstruction'' means gene-inventory reconstruction, not phenotypic.
  \item \textbf{Reference-database version drift.} Both PlasmidFinder (DB 2017-03-19) and VFDB
    have moved on since the paper. Newer database releases could add or reclassify hits.
\end{itemize}

\subsection{Overall reproducibility judgement}

The paper is exemplary in the deposition-and-provenance sense: everything the reader needs to
independently verify structural and functional claims is publicly available and works on the
first try, six years post-publication. The transcriptomics half of the paper (dRNA-seq TSS) is
\textit{depositable} and \textit{re-runnable} but was not attempted here --- so any strong claim
about full replication should be resisted. \textbf{PARTIAL (strong)} is the honest verdict.

\section{Verdict}

\textbf{PARTIAL REPLICATION (strong).}
$\checkmark$ C1, C2, C4, C5, C6, C7, C8, C10 --- 8 of 10 claims independently reproduced.
$\times$ C3 (assembly-step re-run) --- not attempted; assembly product verified instead.
$\times$ C9 (dRNA-seq TSS counts) --- not attempted.

\end{document}
